\chapter{\texttt{medium.atmosphere} --- atmospheric neutrino propagation}
\label{ch:atmosphere}

Atmospheric neutrinos are produced by cosmic-ray air showers at a range of
altitudes above the Earth's surface and propagate down through the
atmosphere -- a comparatively low-density medium, but one whose geometry
(production height, zenith angle, detector depth) is considerably richer
than the single-scale exponential often assumed. This chapter follows the
same profile/geometry/evolutor/probability/flux structure established in
\cref{ch:vacuum}, adapted to that geometry, plus a
production-altitude density layer (\texttt{density.py}) and an atmospheric
slant-depth layer (\texttt{depth.py}) with no direct analogue in the
previous media chapters.

\section{\texttt{profile.py} and \texttt{density.py} --- the atmosphere density model}
\label{sec:atmosphere-profile}

\modulehead{medium/atmosphere/density.py, profile.py}{Pluggable
production-to-surface electron-density backends, and
\pyclass{AtmosphereProfile}, the trajectory container consumed by the
numerical evolutor (\cref{ch:core-numerical}).}

\begin{implbox}
Unlike the Earth chapter's single PREM-derived source
(\cref{sec:earth-profile}), \pyfunc{atmosphere\_density} offers a genuine
choice of density backends, selected by \code{atmosphere\_density\_source}:
\begin{itemize}
  \item \code{"exponential"} (\textbf{the default}): the barometric formula
    $\rho(h)=\rho_0 e^{-h/H}$, with $\rho_0=1.225\times10^{-3}\ \mathrm{g/cm^3}$
    (U.S. Standard Atmosphere sea-level density) and scale height
    $H=8.4\ \mathrm{km}$ -- a single-parameter approximation, adequate when
    only the coarse exponential fall-off matters.
  \item \code{"nusquids"}/\code{"earthatm"}: the external nuSQuIDS
    \code{EarthAtm} density formula (\texttt{external.nusquids}), for
    bit-comparable cross-checks against that reference implementation.
  \item \code{"file"}: linear interpolation of a user-supplied two-column
    altitude/density table.
  \item \code{"mceq"}: densities consistent with an initialised MCEq
    atmosphere object \parencite{Fedynitch2015MCEq}, for internal
    consistency with an MCEq-generated production flux.
  \item \code{"msis"}/\code{"pymsis"}: the NRLMSIS 2.0 empirical
    atmosphere model \parencite{Emmert2021NRLMSIS2}, the most physically
    detailed option (season-, location-, and space-weather-dependent), via
    the optional \texttt{pymsis} backend.
\end{itemize}
Every backend returns mass density $\rho(h)$ in g/cm$^3$; electron density
is obtained via $n_e=Y_e\,\rho\times\text{(mol/g conversion)}$, with the
electron fraction $Y_e$ read from the backend's own configuration (PyMSIS)
or from the shared default $Y_e=0.5$ (air, elsewhere) -- exactly the same
mass-to-electron-density conversion pattern as
\cref{ch:core-common}/\cref{sec:earth-profile}. \pyfunc{atmosphere\_density}
also accepts \code{density\_type="neutron\_density"}, returning
$n_n=(1-Y_e)\,\rho\times\text{(mol/g conversion)}$ from the very same
$Y_e$/$\rho$ already computed for $n_e$ -- no new backend logic, since air's
composition is fully captured by the single scalar $Y_e$ already in use.
This enables the 3+1 sterile neutral-current term (\cref{ch:core-common})
for every density backend uniformly.
\end{implbox}

\begin{implbox}
\pyclass{AtmosphereParameters} bundles the density-backend selection above
with the numerical trajectory settings (\code{nsteps}, sampling
\code{method}) and the analytical-mode fit settings
(\code{perturbative\_segments}, \code{perturbative\_degree}, next section);
\code{matter=False} forces zero density everywhere (a pure-vacuum check).
\pyclass{AtmosphereProfile} builds the production-to-surface
\pyclass{Trajectory} (\cref{ch:core-numerical}) from
(\code{h\_km},\code{theta\_deg},\code{depth\_km}) via the geometry of the
next section, and samples the selected density backend once per segment.
\code{AtmosphereParameters.include\_matter\_nc=False} additionally makes
\pyclass{AtmosphereProfile} sample \code{density\_type="neutron\_density"}
alongside $n_e$, exposed as \code{n\_n\_molcm3} (\code{None} when not
requested); since \pyfunc{atmosphere\_density} supports every backend
uniformly (implbox above), unlike \cref{ch:earth} there is no profile model
for which this can fail to be available.
\end{implbox}

\section{\texttt{geometry.py} and \texttt{depth.py} --- angles, altitude, and atmospheric depth}
\label{sec:atmosphere-geometry}

\modulehead{medium/atmosphere/geometry.py, depth.py}{Zenith/nadir angle
conversions, surface/detector angle mapping, path-length geometry, and the
altitude $\leftrightarrow$ atmospheric-depth relation.}

\begin{notebox}[title={Two Angles and Two ``Depths''}]
This chapter uses \textbf{three} angle labels and \textbf{two} unrelated
quantities both called ``depth'' -- keeping them apart is essential:
\begin{itemize}
  \item $\theta$ (\textbf{detector zenith angle}, the MCEq convention):
    measured at the detector from the local upward vertical;
    $\theta=0^\circ$ is straight down (shortest path), $\theta=90^\circ$ is
    horizontal.
  \item $\eta$ (\textbf{nadir angle}, \cref{ch:earth}'s convention):
    $\eta=\pi-\theta$, identical relation and identical meaning to the one
    used throughout the Earth chapter.
  \item $\alpha$ (\textbf{surface zenith angle}): the zenith angle the same
    ray would have if measured by an observer \emph{at the Earth's
    surface} directly above/along the trajectory, rather than at the
    (possibly buried) detector -- related to $\theta$ by impact-parameter
    conservation (below), not simple equality.
  \item \code{depth\_km} or \code{depth\_m} (\textbf{detector depth}): how far
    below the Earth's surface the detector sits -- a \emph{geometric}
    length, identical in meaning to Earth's $d$ (\cref{ch:earth}).
  \item $X$ (\textbf{atmospheric depth}, \texttt{depth.py}): a
    \emph{column density} in g/cm$^2$ (mass per unit area integrated along
    the path), \textbf{not} a length -- the physical quantity most
    directly tied to how many interaction lengths of atmosphere a
    cosmic-ray shower (and hence its neutrino production point) has
    traversed. It is unrelated to detector depth beyond sharing the English
    word.
\end{itemize}
\end{notebox}

\begin{theorybox}
\subsection*{Angle conversions}
The nadir/zenith relation is exactly the one used throughout
\cref{ch:earth}:
\begin{equation}
  \keyresult{\eta = \pi - \theta}.
\end{equation}
The detector zenith angle $\theta$ and the surface zenith angle $\alpha$ of
the \emph{same} straight-line ray are related by conservation of the
impact parameter (the perpendicular distance from Earth's centre to the
ray is the same whether measured from the detector's radius $r_d=R_\oplus-d$
or the surface radius $R_\oplus$):
\begin{equation}
  \keyresult{r_d\sin\theta = R_\oplus\sin\alpha}
  \qquad\Longleftrightarrow\qquad
  \theta=\arcsin\!\Bigl(\frac{R_\oplus}{r_d}\sin\alpha\Bigr),\
  \alpha=\arcsin\!\Bigl(\frac{r_d}{R_\oplus}\sin\theta\Bigr).
\end{equation}
Because $R_\oplus/r_d>1$ for a buried detector, not every surface angle
maps to a valid detector ray: $\alpha$ must stay below
$\alpha_{\max}=\arcsin(r_d/R_\oplus)$ for the first arcsine to remain
defined (\pyfunc{alpha\_max\_for\_detector\_depth}). For a surface detector
($d=0$), $r_d=R_\oplus$ and $\alpha\equiv\theta$ trivially.

\subsection*{Path lengths}
For a neutrino produced at altitude $h$ and detected at zenith angle
$\theta$ with detector depth $d$ (so $r_d=R_\oplus-d$), the straight-line
law of cosines applied to the triangle (Earth centre, detector, production
point) gives the total detector-to-production distance,
\begin{equation}
  \keyresult{L_{\rm total} = -r_d\cos\theta+\sqrt{(R_\oplus+h)^2-r_d^2\sin^2\theta}},
\end{equation}
and, setting $h=0$, the purely \textbf{underground} part of the same ray
(detector to surface),
\begin{equation}
  L_{\rm underground} = -r_d\cos\theta+\sqrt{R_\oplus^2-r_d^2\sin^2\theta}.
\end{equation}
The \textbf{atmospheric} path length actually carrying matter effects is
the difference, $L_{\rm atm}=L_{\rm total}-L_{\rm underground}$
(\pyfunc{atmosphere\_path\_length}); \pyfunc{atmosphere\_path\_grid} samples
altitude along it via the same law-of-cosines relation evaluated at
distance $s$ from the detector,
\begin{equation}
  \keyresult{r(s)=\sqrt{r_d^2+s^2+2r_ds\cos\theta}},\qquad h(s)=r(s)-R_\oplus.
\end{equation}

\subsection*{Vertical and slant atmospheric depth}
Independently of the path-length geometry above, the atmospheric depth $X$
(the column density integral, used e.g.\ to tabulate cosmic-ray shower
development or MCEq/Honda fluxes rather than oscillation propagation
directly) is
\begin{equation}
  \keyresult{X_{\rm vertical}(h) = \int_h^\infty \rho(h')\,\mathrm{d}h'},
  \qquad
  X_{\rm slant}(h,\alpha) = \frac{X_{\rm vertical}(h)}{\cos\alpha},
\end{equation}
the second equation being the plane-parallel approximation (valid for
$\alpha$ not too close to the horizon, where the Earth's curvature would
need to be accounted for); $\alpha$ here is the \emph{surface} zenith
angle defined above, since the plane-parallel picture is anchored to a
flat patch of atmosphere.
\end{theorybox}

\begin{implbox}
\begin{itemize}
  \item \pyfunc{theta\_to\_eta} / \pyfunc{eta\_to\_theta}: $\eta=\pi-\theta$
    and its inverse.
  \item \pyfunc{alpha\_surface\_to\_theta\_detector} /
    \pyfunc{theta\_detector\_to\_alpha\_surface} /
    \pyfunc{alpha\_max\_for\_detector\_depth}: the impact-parameter mapping
    above and its domain limit.
  \item \pyfunc{total\_path\_length} / \pyfunc{underground\_path\_length} /
    \pyfunc{atmosphere\_path\_length}: the three path-length formulas
    above.
  \item \pyfunc{altitude\_along\_detector\_path} /
    \pyfunc{atmosphere\_path\_grid}: the $r(s)$ sampling formula above,
    and the resulting distance/altitude grid used to build an
    \pyclass{AtmosphereProfile} trajectory.
  \item \pyfunc{atmosphere\_vertical\_depth(h\_km, rho\_gcm3)} /
    \pyfunc{atmosphere\_slant\_depth(X\_vertical, alpha\_deg)}
    (\texttt{depth.py}): the $X_{\rm vertical}$/$X_{\rm slant}$ formulas
    above, the first via trapezoidal integration from the top of the
    supplied grid downward.
  \item \pyfunc{compute\_dXdh} / \pyfunc{interpolate\_flux\_at\_Xobs}
    (\texttt{depth.py}): the depth-profile derivative and a
    depth-tabulated-flux interpolator, both backend-neutral utilities for
    working with $\Phi(E,X,\alpha)$ tables from any external source.
\end{itemize}
\end{implbox}

\section{\texttt{evolutor.py} --- atmosphere evolution operators}
\label{sec:atmosphere-evolutor}

\modulehead{medium/atmosphere/evolutor.py}{\pyfunc{atmosphere\_evolutor}:
dispatches between a numerical (default) and an analytical perturbative
atmosphere evolutor.}

\begin{implbox}
The same numerical/analytical distinction of \cref{ch:earth}
applies here, with the \textbf{default reversed} (\code{method="numerical"}
by default, versus Earth's \code{"analytical"} default): the atmosphere's
low density and short path make the numerical matrix-exponential pipeline
(\cref{ch:core-numerical}) cheap enough to use by default.
\begin{itemize}
  \item \pyfunc{atmosphere\_evolutor\_numerical(...)}: builds an
    \pyclass{AtmosphereProfile} trajectory (\cref{sec:atmosphere-profile})
    and calls \pyfunc{evolutor\_numerical} (\cref{ch:core-numerical})
    directly on its sampled density; entries with $L_{\rm atm}\le0$
    (grazing/no-path geometries) are masked to the identity.
  \item \pyfunc{atmosphere\_evolutor\_analytical(...)}: \emph{fits} a local
    polynomial density model on the fly -- splits the atmospheric path
    into \code{perturbative\_segments} equal sub-intervals, samples the
    density backend at \code{perturbative\_degree}$+1$ Chebyshev-like nodes
    per sub-interval (\code{q}$\in[-1,1]$, centred/rescaled), and builds an
    \pyclass{AtmospherePolynomialProfile}
    (\cref{ch:core-perturbative}, \texttt{atmosphere} model) from those
    samples, then evolves it with
    \pyfunc{evolutor\_perturbative\_segment} (\cref{ch:core-perturbative})
    exactly as the Earth analytical pipeline does. Unlike Earth's
    tabulated shell coefficients, the polynomial here is fitted
    \emph{per call}, from whichever density backend is selected
    (\cref{sec:atmosphere-profile}) -- the price of supporting arbitrary
    density sources with the same closed-form perturbative machinery.
  \item \pyfunc{atmosphere\_evolutor(...)}: the public dispatcher, mirroring
    \pyfunc{earth\_probability\_state} (\cref{sec:earth-probability}) in
    spirit.
\end{itemize}
\end{implbox}

\begin{implbox}
Both evolutors accept \code{atmosphere.include\_matter\_nc=False}
(\cref{sec:atmosphere-profile}), enabling the 3+1 sterile neutral-current
term (\cref{ch:core-common}). The numerical evolutor forwards
\pyclass{AtmosphereProfile}'s sampled \code{n\_n\_molcm3} straight through
to \pyfunc{evolutor\_numerical} as \code{n\_n\_mol\_cm3}. The analytical
evolutor fits a \emph{second}, independent
\pyclass{AtmospherePolynomialProfile} to a neutron-density sample at the
same nodes as the electron-density fit, and passes its coefficients as
\code{coefficients\_n} into the same segment model
(\cref{ch:core-perturbative}), then calls
\pyfunc{evolutor\_perturbative\_segment} with
\code{include\_matter\_nc=True}. Both default to \code{False}, reproducing
the pre-existing CC-only behaviour exactly.
\end{implbox}

\section{\texttt{probability.py} --- atmosphere-surface probabilities}
\label{sec:atmosphere-probability}

{\sloppy\modulehead{medium/atmosphere/probability.py}{\pyfunc{atmosphere\_\allowbreak{}probability\_\allowbreak{}transition}
and \pyfunc{atmosphere\_\allowbreak{}probability\_\allowbreak{}state}: final flavour probabilities at
the Earth surface from an atmosphere evolution operator.
\pyfunc{atmosphere\_\allowbreak{}probability\_\allowbreak{}integrated},
\pyfunc{atmosphere\_\allowbreak{}probability\_\allowbreak{}integrated\_\allowbreak{}angular}, and
\pyfunc{atmosphere\_\allowbreak{}probability\_\allowbreak{}integrated\_\allowbreak{}height} build on
\pyfunc{atmosphere\_\allowbreak{}probability\_\allowbreak{}state} to average over energy, the full
zenith range, and production height, respectively.}}

\begin{implbox}
Structurally identical to \cref{ch:vacuum,ch:earth}:
\pyfunc{atmosphere\_\allowbreak{}probability\_\allowbreak{}transition(oscillation, E\_MeV, h\_km, theta\_deg, depth\_km, method, ...)}
returns the full transition-probability matrix $P=|S_{\rm atm}|^2$
(\cref{ch:core-common}); \pyfunc{atmosphere\_\allowbreak{}probability\_\allowbreak{}state(nustate, ..., massbasis, ...)}
dispatches between the coherent formula
($\psi_{\rm surface}=S_{\rm atm}\psi_{\rm initial}$,
$P_\alpha=|\psi_{\rm surface,\alpha}|^2$) and the incoherent formula
($P_\alpha=\sum_i|(S_{\rm atm}U_{\rm PMNS})_{\alpha i}|^2w_i$) of
\cref{ch:core-common}, exactly as \pyfunc{vacuum\_probability\_state} and
\pyfunc{earth\_probability\_state} do,
using whichever evolutor \pyfunc{atmosphere\_evolutor}
(\cref{sec:atmosphere-evolutor}) returns. Because every function in this
section takes the whole \pyclass{AtmosphereParameters} bundle as its
\code{atmosphere} argument rather than exploding it into individual
keywords, \code{include\_matter\_nc} (previous section) is already
available at every level here with no further plumbing.
\end{implbox}

\begin{notebox}
\textbf{A real, previously uncaught 3-flavour-only bug.}
\pyfunc{atmosphere\_probability\_state} used to crash outright for a
4-flavour (3+1 sterile) initial state, even though
\pyfunc{atmosphere\_evolutor} underneath it already built the correct
$4\times4$ operator: a shared broadcasting helper
(\code{util.type.broadcast\_last3}, used identically in
\cref{ch:vacuum}) hardcoded its input's final dimension to exactly 3. The
fix generalises the helper to accept 3 \emph{or} 4 (renamed
\code{broadcast\_flavour\_vector} to match, since it is no longer
3-specific) and is shared by \texttt{medium.vacuum.probability},
\texttt{medium.vacuum.evolutor}, and this module. No test exercised a
sterile initial state through this function before the fix.
\end{notebox}

\begin{implbox}
Atmospheric probabilities depend on three axes at once ($E$, production
height $h$, zenith angle $\theta$), so this medium is the only one with
three dedicated \code{*\_integrated} reductions, each delegating to the
corresponding \cref{ch:core-common} primitive:
\begin{itemize}
  \item \pyfunc{atmosphere\_probability\_integrated(nustate, ..., spectrum, integrate\_angular, integrate\_height, ...)}:
    \textbf{always} averages over energy (spectrum-weighted, via
    \pyfunc{core.common.probability.probability\_integrated}), then
    optionally chains a height average and/or an angular average by hand,
    each a further, independent reduction rather than one fused operation.
  \item \pyfunc{atmosphere\_probability\_integrated\_angular(nustate, ...)}:
    averages over the full zenith range via
    \pyfunc{core.common.probability.probability\_integrated\_angular}
    (geometric $\sin\theta$ weight).
  \item \pyfunc{atmosphere\_probability\_integrated\_height(nustate, ..., production\_flux, ...)}:
    averages over production height, weighted by the height-resolved
    production flux, again via
    \pyfunc{core.common.probability.probability\_integrated} -- height is
    not a universal \cref{ch:core-common} concept, so this wrapper lives
    here, even though it reuses that module's generic weighted-average
    primitive.
\end{itemize}
Every reduction's \code{*\_dim} parameter defaults to $-2$ (``the axis
immediately before flavour''), matching \cref{ch:core-common}'s own
convention; for the composite function's fixed energy-then-height-then-angle
chaining order to line up with that default at every step, the input grid
should be laid out as $(\ldots,\theta,h,E,\text{flavour})$.
\end{implbox}

\section{\texttt{flux.py} --- atmosphere flux}
\label{sec:atmosphere-flux}

{\sloppy\modulehead{medium/atmosphere/flux.py}{\pyfunc{atmosphere\_\allowbreak{}flux\_\allowbreak{}state}
(surface flux from \pyfunc{atmosphere\_\allowbreak{}probability\_\allowbreak{}state}) and three
\code{*\_integrated} variants mirroring the probability-layer reductions
above.}}

\begin{implbox}
Like \texttt{probability.py} above, every function here takes the whole
\code{atmosphere} (\pyclass{AtmosphereParameters}) bundle, so
\code{include\_matter\_nc} is available with no further plumbing.
\begin{itemize}
  \item \pyfunc{atmosphere\_flux\_state(nustate, oscillation, E\_MeV, h\_km, theta\_deg, flux, spectrum, depth\_km, ...)}:
    the standard pattern -- \pyfunc{atmosphere\_probability\_state} then
    \pyfunc{core.common.flux.flux\_state} (\cref{ch:core-common}).
  \item \pyfunc{atmosphere\_flux\_integrated(..., integrate\_angular, integrate\_height, ...)}:
    always integrates over energy via
    \pyfunc{core.common.flux.flux\_integrated} (an unnormalised rate,
    unlike the probability-layer average), then optionally chains a plain
    trapezoidal height integral and/or an angular integral
    (\pyfunc{core.common.flux.flux\_integrated\_angular}) -- no separate
    height weight is needed here, since \code{flux} already carries
    whatever height dependence it has.
  \item \pyfunc{atmosphere\_flux\_integrated\_angular(...)}: integrates over
    the full zenith range via
    \pyfunc{core.common.flux.flux\_integrated\_angular}.
  \item \pyfunc{atmosphere\_flux\_integrated\_height(..., production\_flux, ...)}:
    integrates over production height using the tabulated production-flux
    data. Atmosphere-specific and, unlike the three functions above, calls
    no \texttt{core.common} function directly: it combines
    \pyfunc{atmosphere\_probability\_integrated\_height} (the
    production-flux-weighted average probability over height, $\langle
    P\rangle_h$) with the plain trapezoidal integral of
    \code{production\_flux} itself over height,
    $\Phi_h=\int\text{production\_flux}(h)\,\mathrm{d}h$; the product
    $\langle P\rangle_h\,\Phi_h$ equals $\int P(h)\,\text{production\_flux}(h)\,\mathrm{d}h$
    exactly, by the definition of the weighted average.
\end{itemize}
\end{implbox}

\begin{notebox}[title={Detector Composition Lives in the Pipeline Layer}]
Unlike \cref{ch:earth}, \texttt{medium.atmosphere} only ever propagates
production $\to$ surface: it has no notion of a detector, and no
\texttt{validation.py}. Weighting a surface or detector probability by
production-flux data, integrating a detector flux over production height,
and summing the incoherent contributions from several produced flavours are
all pipeline-level concerns, composed in
\texttt{pipeline.atmosphere\_earth} (\pyfunc{detector\_flux\_from\_production},
\pyfunc{integrate\_detector\_flux\_over\_height},
\pyfunc{sum\_detected\_flavours}) on top of
\pyfunc{pipeline.atmosphere\_earth.propagate\_surface\_to\_detector}, which
applies the Earth evolutor of \cref{ch:earth} to this medium's coherent
surface state via ordinary matrix-vector composition, not the matrix-matrix
composition used internally by \cref{sec:earth-evolutor}.
\end{notebox}
